% Values: PRECISE, FAITHFUL
% Evidence planned:
% - PRECISE: one section per abstract parameter or unspecified point of the
%   body (IntRange and norm, characters, Float operations, equality, bitwise
%   and string conversion, evaluation order and Debug, resource limits,
%   platform and ports, eco extensions); each states what the body requires,
%   with a reference to the body, then a table with one row per
%   implementation (Elm 0.19.1 on JavaScript; eco 0.2.0) giving its choice and
%   the consequences for programs.
% - FAITHFUL: every ruling of the ambiguity log marked UNSPECIFIED or CUSTOM
%   that names this appendix, and every "recorded in the appendix" ruling
%   (A8, A9, A13, A14, A33, A34, A39, A40 to A55, A57 to A59, A61 to
%   A72, A76, A79 to A81, A83, A92), cited beside the text it shaped; every
%   implementation fact traceable to the survey ("Implementation-Dependent
%   Points", "Arithmetic", "Equality", "Comparison and Order", "Char and String
%   Basics", "Bitwise", "Remaining Runtime Failures", "Ports at Runtime", "Eco
%   Differences") or to the ruling; a closing table of known divergences with
%   the ruling that classifies each.

\chapter{Implementation-Dependent Details}\label{ch:appendix-implementation}

Where implementations of Standard Elm legitimately differ, the body of the
definition states only what every implementation must satisfy, using an
abstract parameter where one is needed. % D8
This appendix gives, for each such parameter or unspecified point, what the
body requires of it and what two implementations choose. The implementations
are Elm 0.19.1 with the elm/core 1.0.5 kernel on JavaScript, written
\emph{Elm/JS}, and eco 0.2.0 producing native code, written \emph{eco}; eco's
output for JavaScript behaves as Elm/JS does, except for \kw{Bytes} at the
boundary (Section~\ref{sec:appendix-implementation-bytes}). Each section states the
requirement, with a reference to the body, and then gives a table with one row
per implementation: its choice, and the consequences for programs. An
implementation \emph{conforms} on a point when its choice meets the
requirement. Section~\ref{sec:appendix-implementation-divergences} collects
every known point on which an implementation does not conform, with the
ruling that settles it. The facts recorded here describe the versions named;
they are not part of the definition.

% ===========================================================================
\section{The Int Range and Normalisation}\label{sec:appendix-implementation-int}

The body requires the following of the type \kw{Int}
(Chapter~\ref{ch:syntax-core}, Section~\ref{sec:dynamic-core-simple} and
Section~\ref{sec:appendix-dynamic-basis-apply}).
\begin{itemize}
\item $\IntRange$ is a set of integers containing every integer $n$ with
  $-2^{31}\leq n\leq 2^{31}-1$; Int has one zero. % A40, A45
\item $\norm$ maps every integer into $\IntRange$ and is the identity on
  $\IntRange$. Int addition, subtraction, multiplication, negation, absolute
  value and power with a non-negative exponent give $\norm$ of the exact
  result, and an integer constant spelling $n$ denotes $\norm(n)$. % A9, A40
  % DECIDED: norm is the identity on IntRange; A40 calls it a normalisation
  % into the range, which is read as leaving members of the range unchanged.
\item For $b\neq 0$, \kw{a // b} is $\norm$ of the truncation of $a/b$ toward
  zero; \kw{x // 0}, \kw{modBy 0 x} and \kw{remainderBy 0 x} are $0$.
  % A41, A42, D3
\item For $e<0$, \kw{b \textasciicircum{} e} at Int is $1$ when $b=1$, $1$ or
  $-1$ by the parity of $e$ when $b=-1$, and $0$ otherwise. % A43
\item \kw{round}, \kw{floor}, \kw{ceiling} and \kw{truncate} give the integer
  they denote when it lies in $\IntRange$, the least or greatest member of
  $\IntRange$ when it lies below or above it (an infinity included), and $0$
  for NaN. % A44
\end{itemize}
A value of type \kw{Int} is always a member of $\IntRange$: never a fraction,
NaN, an infinity or a negative zero.

\subsection{The Range and Normalisation}
% A9, A40
\begin{center}\small
\begin{tabular}{@{}>{\raggedright\arraybackslash}p{0.10\textwidth}>{\raggedright\arraybackslash}p{0.43\textwidth}>{\raggedright\arraybackslash}p{0.39\textwidth}@{}}
\hline
Impl. & Choice & Consequences for programs\\
\hline
Elm/JS & An Int is an IEEE 754 double. \kw{+}, \kw{-} and \kw{*} give the exact
result rounded to the nearest double: exact while its magnitude is at most
$2^{53}$, rounded above that, an infinity beyond the largest double. An
integer constant is reduced modulo $2^{64}$ into $[-2^{63},2^{63})$ as two's
complement. &
No $\IntRange$ and $\norm$ describe Elm/JS: Int results above $2^{53}$ in
magnitude are inexact, an overflow gives an infinity, and other operations
below leave the integers. Programs whose Int results stay within $\pm 2^{53}$
compute exactly.\\
eco & $\IntRange$ is $[-2^{63},2^{63}-1]$, and $\norm(n)$ is the $m$ in
$\IntRange$ with $m\equiv n \pmod{2^{64}}$ (64-bit two's complement wrap).
Integer constants likewise: \kw{0xFFFFFFFFFFFFFFFF} denotes $-1$. &
Conforms. Overflow wraps silently: the greatest Int plus $1$ is the least;
\kw{abs} and \kw{negate} of the least Int give it back.\\
\hline
\end{tabular}
\end{center}

\subsection{Integer Division and Remainders}
% A41, A42
\begin{center}\small
\begin{tabular}{@{}>{\raggedright\arraybackslash}p{0.10\textwidth}>{\raggedright\arraybackslash}p{0.43\textwidth}>{\raggedright\arraybackslash}p{0.39\textwidth}@{}}
\hline
Impl. & Choice & Consequences for programs\\
\hline
Elm/JS & \kw{a // b} is the quotient $a/b$ converted to a 32-bit integer
(truncation, then reduction modulo $2^{32}$ into $[-2^{31},2^{31})$; NaN and
the infinities give $0$), so \kw{x // 0} is $0$. \kw{modBy 0 x} stops the
program with the message ``Cannot perform mod 0. Division by zero error.''
\kw{remainderBy 0 x} is NaN, held as an Int. &
Quotients outside $\pm 2^{31}$ wrap: \kw{5000000000 // 1} is
\kw{705032704}. \kw{modBy 0} ends the program; NaN from \kw{remainderBy 0}
spreads through later arithmetic (\kw{remainderBy 0 7 + 1} is NaN).\\
eco & Truncated division on 64 bits; \kw{x // 0}, \kw{modBy 0 x} and
\kw{remainderBy 0 x} are $0$; the least Int \kw{// -1} is the least Int. &
Conforms.\\
\hline
\end{tabular}
\end{center}

\subsection{Powers with a Negative Exponent}
% A43
\begin{center}\small
\begin{tabular}{@{}>{\raggedright\arraybackslash}p{0.10\textwidth}>{\raggedright\arraybackslash}p{0.43\textwidth}>{\raggedright\arraybackslash}p{0.39\textwidth}@{}}
\hline
Impl. & Choice & Consequences for programs\\
\hline
Elm/JS & The real power, held as an Int: \kw{2 \textasciicircum{} -1} is
\kw{0.5} and \kw{0 \textasciicircum{} -1} is an infinity. &
Does not conform: the result is not an integer.\\
eco & $0$ for every base, \kw{1} and \kw{-1} included. For $e\geq 0$, the
power reduced modulo $2^{64}$, as $\norm$. &
Does not conform for the bases $1$ and $-1$: \kw{1 \textasciicircum{} -1} is
$0$, not $1$.\\
\hline
\end{tabular}
\end{center}

\subsection{Conversion from Float}
% A44
\begin{center}\small
\begin{tabular}{@{}>{\raggedright\arraybackslash}p{0.10\textwidth}>{\raggedright\arraybackslash}p{0.43\textwidth}>{\raggedright\arraybackslash}p{0.39\textwidth}@{}}
\hline
Impl. & Choice & Consequences for programs\\
\hline
Elm/JS & \kw{round}, \kw{floor} and \kw{ceiling} of NaN or an infinity give NaN
or that infinity, held as an Int. \kw{truncate} converts to a 32-bit integer
(NaN and the infinities give $0$). &
Does not conform: \kw{round (0/0)} is NaN, and \kw{truncate 3e10} is
\kw{-64771072}.\\
eco & Saturating conversion to 64 bits: NaN gives $0$; values above or below
$\IntRange$, infinities included, give its greatest or least member. &
Conforms: \kw{truncate 3e10} is \kw{30000000000}, and \kw{round (1/0)} and
\kw{truncate 1e20} are \kw{9223372036854775807}.\\
\hline
\end{tabular}
\end{center}

\subsection{Negative Zero}
% A45
\begin{center}\small
\begin{tabular}{@{}>{\raggedright\arraybackslash}p{0.10\textwidth}>{\raggedright\arraybackslash}p{0.43\textwidth}>{\raggedright\arraybackslash}p{0.39\textwidth}@{}}
\hline
Impl. & Choice & Consequences for programs\\
\hline
Elm/JS & \kw{negate 0}, \kw{0 * -1}, \kw{modBy 4 -4}, \kw{remainderBy 4 -4},
\kw{round -0.4} and \kw{ceiling -0.5} give a negative zero held as an Int. &
Does not conform. The sign shows after conversion:
\kw{1 / toFloat (negate 0)} is \kw{-Infinity}. Equality does not see it.\\
eco & Int has one zero. & Conforms: \kw{1 / toFloat (negate 0)} is
\kw{Infinity}.\\
\hline
\end{tabular}
\end{center}

In Elm/JS a value of type \kw{Int} may thus be a fraction, NaN, an infinity or
a negative zero, none of which is a member of any $\IntRange$; Elm/JS does not
conform at the extremes of Int. A program whose exact Int results, quotients
and conversions from Float all lie between $-2^{31}$ and $2^{31}-1$, and which
applies neither \kw{modBy 0}, \kw{remainderBy 0} nor Int
\kw{\textasciicircum} with a negative exponent, has the same meaning in every
implementation listed here, apart from the sign of a zero seen through
\kw{toFloat}.

% ===========================================================================
\section{Characters and Strings}\label{sec:appendix-implementation-chars}

The body requires the following (Chapter~\ref{ch:syntax-core},
Section~\ref{sec:dynamic-core-simple} and
Section~\ref{sec:appendix-dynamic-basis-apply}).
\begin{itemize}
\item \kw{Char} is an abstract class of characters with an injective function
  \kw{toCode} into the integers; a string is a finite sequence of
  characters. % A52
\item The escape \kw{\textbackslash u\kwlb}$h$\kw{\kwrb} denotes a Unicode
  scalar value: surrogate code points are excluded. % A8
\item Chars are ordered by \kw{toCode}; strings are ordered lexicographically
  by their sequences of characters. % A53
\item The letters of identifiers are classified by the general categories of
  a version of the Unicode standard chosen by the implementation. % A14
\item Case mapping is outside the initial basis and has no requirement. % A54
\end{itemize}

\subsection{What a Character Is}
% A52, A8
\begin{center}\small
\begin{tabular}{@{}>{\raggedright\arraybackslash}p{0.10\textwidth}>{\raggedright\arraybackslash}p{0.43\textwidth}>{\raggedright\arraybackslash}p{0.39\textwidth}@{}}
\hline
Impl. & Choice & Consequences for programs\\
\hline
Elm/JS & A Char is a JavaScript string of one code point, one or two UTF-16
code units; case mapping can make it longer. \kw{toCode} gives its first code
point, and NaN for a lone high surrogate. \kw{fromCode n} gives U+FFFD for
$n<0$ or $n>\mathrm{10FFFF}_{16}$, and passes surrogates through. A string is
a sequence of UTF-16 code units. A surrogate escape is accepted. &
\kw{toCode} is not injective on the Chars Elm/JS can make: the two-character
result of \kw{Char.toUpper \textquotesingle\ss\textquotesingle} has the code
of \kw{\textquotesingle S\textquotesingle}. A string holding a lone surrogate
makes \kw{String.toList} run until memory is exhausted.\\
eco & A Char is one code point in $0$ to $\mathrm{10FFFF}_{16}$; a surrogate
code point is a Char of its own. \kw{toCode} is the code point; \kw{fromCode}
as Elm/JS. A string is a sequence of UTF-16 code units: a valid surrogate pair
is one character, a lone surrogate is one character. A surrogate escape is
accepted. &
\kw{toCode} is injective. \kw{String.fromList (String.toList s)} is \kw{s}
for every string.\\
\hline
\end{tabular}
\end{center}
In both, \kw{String.length} and string indices count UTF-16 code units, so a
character above U+FFFF counts twice: the length of a three-character string
whose middle character is U+1F600 is $4$.

\subsection{The Order of Characters and Strings}
% A53
\begin{center}\small
\begin{tabular}{@{}>{\raggedright\arraybackslash}p{0.10\textwidth}>{\raggedright\arraybackslash}p{0.43\textwidth}>{\raggedright\arraybackslash}p{0.39\textwidth}@{}}
\hline
Impl. & Choice & Consequences for programs\\
\hline
Elm/JS & Chars and strings both by UTF-16 code units. &
The two orders agree with each other but not with \kw{toCode}: a character
above U+FFFF sorts below the characters U+E000 to U+FFFF. Does not
conform.\\
eco & Chars by code point; strings by UTF-16 code units. &
Char order conforms; string order does not, for strings comparing a character
above U+FFFF with one in U+E000 to U+FFFF. Char order and the order of
one-character strings then disagree.\\
\hline
\end{tabular}
\end{center}

\subsection{The Unicode Version of Identifiers}
% A14
\begin{center}\small
\begin{tabular}{@{}>{\raggedright\arraybackslash}p{0.10\textwidth}>{\raggedright\arraybackslash}p{0.43\textwidth}>{\raggedright\arraybackslash}p{0.39\textwidth}@{}}
\hline
Impl. & Choice & Consequences for programs\\
\hline
Elm/JS & The Unicode tables of the Haskell system that built the compiler. &
A letter assigned in a later Unicode version than one implementation knows is
not a letter there: an identifier using it is accepted by one implementation
and rejected by the other.\\
eco & Unicode 14.0.0. & \\
\hline
\end{tabular}
\end{center}

\subsection{Case Mapping}
% A54
\begin{center}\small
\begin{tabular}{@{}>{\raggedright\arraybackslash}p{0.10\textwidth}>{\raggedright\arraybackslash}p{0.43\textwidth}>{\raggedright\arraybackslash}p{0.39\textwidth}@{}}
\hline
Impl. & Choice & Consequences for programs\\
\hline
Elm/JS & Full case mapping of the host's Unicode version, with the final sigma
rule; \kw{toLocaleUpper} and \kw{toLocaleLower} use the host's locale in a
browser. &
\kw{Char.toUpper \textquotesingle\ss\textquotesingle} is the two-character
\kw{"SS"}, held as a Char. \kw{String.toLower} lowers a word-final capital
sigma to final sigma.\\
eco & Simple case mapping of Unicode 17.0, one code point to one; no final
sigma rule; \kw{toLocaleUpper} and \kw{toLocaleLower} are \kw{toUpper} and
\kw{toLower}. &
\kw{Char.toUpper \textquotesingle\ss\textquotesingle} is
\kw{\textquotesingle\ss\textquotesingle}; case mapping keeps the length of a
string.\\
\hline
\end{tabular}
\end{center}
The predicates of the module \kw{Char} (\kw{isUpper}, \kw{isLower},
\kw{isAlpha}, \kw{isDigit}, \kw{isOctDigit} and \kw{isHexDigit}) are, in
both, defined by their elm/core text over ASCII code ranges, and agree.

% ===========================================================================
\section{Floating-Point Operations}\label{sec:appendix-implementation-float}

The body requires the following of the type \kw{Float}
(Section~\ref{sec:dynamic-core-simple} and
Section~\ref{sec:appendix-dynamic-basis-apply}).
\begin{itemize}
\item \kw{compare} on Floats is a total order: $-0$ and $+0$ are equal, NaN
  equals NaN and is greater than every other Float. % A46
\item At Float, \kw{<}, \kw{<=}, \kw{>} and \kw{>=} are the IEEE predicates,
  false when an operand is NaN; on tuples and lists they are the
  lexicographic order derived from \kw{compare}. % A47
\item At Float, \kw{min} and \kw{max} are the \emph{minimum} and
  \emph{maximum} of IEEE 754-2019: NaN when either argument is NaN, and
  $-0$ below $+0$, whatever the order of the arguments. At every other
  comparable type they give the results of their elm/core definitions. % A48
\item At Float, \kw{clamp lo hi x} is \kw{min (max x lo) hi}, with the
  \kw{min} and \kw{max} above: $hi$ when $lo>hi$, NaN when any argument is
  NaN, and \kw{clamp 0 1 -0} is $+0$. At Int it is given by its elm/core
  definition through \kw{<} and \kw{>}: $lo$ when $x<lo$, else $hi$ when
  $x>hi$, else $x$, so that \kw{clamp 10 0 5} is $10$. % A48, A49, A60
\item At Float, \kw{+}, \kw{-}, \kw{*}, \kw{/}, \kw{sqrt} and \kw{negate}
  are the operations of IEEE 754-2019, correctly rounded, with no latitude
  in accuracy (Figure~\ref{fig:dynamic-basis-float}).
\item The power and the transcendental functions (the trigonometric functions,
  their inverses, \kw{atan2} and the logarithm under \kw{logBase}) denote the
  real function rounded to a Float, with special values as IEEE 754-2019
  recommends; their accuracy is implementation-dependent. % A65
\end{itemize}
In both implementations a Float is an IEEE 754 binary64 number, and \kw{+},
\kw{-}, \kw{*}, \kw{/} and \kw{sqrt} are the IEEE operations rounding to
nearest.
% Audit Y9: IEEE 754 requires a correctly rounded square root, so sqrt is
% with + - * / (Appendix D); A65 covers only the power and the
% transcendental functions.

\subsection{\kw{compare} with NaN}
% A46
\begin{center}\small
\begin{tabular}{@{}>{\raggedright\arraybackslash}p{0.10\textwidth}>{\raggedright\arraybackslash}p{0.43\textwidth}>{\raggedright\arraybackslash}p{0.39\textwidth}@{}}
\hline
Impl. & Choice & Consequences for programs\\
\hline
Elm/JS & \kw{LT} if $a<b$, \kw{EQ} if $a$ and $b$ are equal, else \kw{GT}: any
comparison with NaN is \kw{GT}. &
Does not conform: \kw{compare 1 nan} and \kw{compare nan 1} are both
\kw{GT}. Sorting, \kw{Dict} and \kw{Set} with NaN keys have no defined
result.\\
eco & The total order of the body. &
Conforms: \kw{compare 1 nan} is \kw{LT} and \kw{compare nan nan} is \kw{EQ};
NaNs sort last and are one key of a \kw{Dict} or \kw{Set}.\\
\hline
\end{tabular}
\end{center}

\subsection{The Order Operators}
% A47
\begin{center}\small
\begin{tabular}{@{}>{\raggedright\arraybackslash}p{0.10\textwidth}>{\raggedright\arraybackslash}p{0.43\textwidth}>{\raggedright\arraybackslash}p{0.39\textwidth}@{}}
\hline
Impl. & Choice & Consequences for programs\\
\hline
Elm/JS & The IEEE predicate when an operand is written as a constant;
otherwise derived from \kw{compare} above. &
Does not conform: with \kw{one = 1}, \kw{one > nan} is \kw{True} while
\kw{1 > nan} is \kw{False}, and \kw{nan >= nan} is \kw{True} when neither
operand is a constant.\\
eco & IEEE predicates at Float on every path; on tuples and lists, by
\kw{compare}. &
Conforms: \kw{nan < 1} is \kw{False}, and \kw{[ nan ] > [ 1 ]} is
\kw{True}.\\
\hline
\end{tabular}
\end{center}

\subsection{\kw{min}, \kw{max} and \kw{clamp}}
% A48, A49
\begin{center}\small
\begin{tabular}{@{}>{\raggedright\arraybackslash}p{0.10\textwidth}>{\raggedright\arraybackslash}p{0.43\textwidth}>{\raggedright\arraybackslash}p{0.39\textwidth}@{}}
\hline
Impl. & Choice & Consequences for programs\\
\hline
Elm/JS & The elm/core definitions, at every type, through the comparison of
\kw{compare} above, which makes any comparison with NaN \kw{GT} and the two
zeros \kw{EQ}. &
\kw{min} and \kw{max} do not conform at Float when an argument is NaN or
both are zeros: the result depends on the order of the arguments.
\kw{min nan 1} and \kw{max 1 nan} are \kw{1}, where the body gives NaN.
\kw{clamp} conforms at Int; at Float it does not conform when $lo>hi$, when
an argument is NaN, or with zeros of different signs: \kw{clamp 10 0 5} is
\kw{10}, \kw{clamp 0 1 nan} is \kw{1}, \kw{clamp nan 1 0.5} is \kw{0.5}
and \kw{clamp 0 1 -0} is $-0$, where the body gives $+0$, NaN, NaN and
$+0$.\\
eco & At Float, IEEE 754-2019 \emph{minimum} and \emph{maximum} (NaN in
either argument gives NaN; $-0<+0$), and \kw{clamp lo hi x} is
\kw{min (max x lo) hi}. At every other type, the elm/core definitions. &
\kw{min}, \kw{max} and \kw{clamp} conform.\\
\hline
\end{tabular}
\end{center}

\subsection{Special Values and Accuracy}
% A65
\begin{center}\small
\begin{tabular}{@{}>{\raggedright\arraybackslash}p{0.10\textwidth}>{\raggedright\arraybackslash}p{0.43\textwidth}>{\raggedright\arraybackslash}p{0.39\textwidth}@{}}
\hline
Impl. & Choice & Consequences for programs\\
\hline
Elm/JS & The host's ECMAScript \kw{Math} functions, whose accuracy the host
chooses. \kw{\textasciicircum} at Float follows ECMAScript: \kw{1
\textasciicircum{} (1/0)} is NaN. &
Does not conform for the special values of \kw{\textasciicircum}: IEEE
754-2019 gives $1$. \kw{logBase 10 1000} is \kw{2.9999999999999996}.\\
eco & \kw{\textasciicircum} at Float follows C99 Annex F (\kw{1
\textasciicircum{} (1/0)} is $1$); \kw{sin} and \kw{cos} are those of the
code generator, \kw{tan x} is computed as \kw{sin x / cos x}, and the inverse
trigonometric functions are those of the C mathematics library. &
Accuracy is that of the functions named. The special values of
\kw{\textasciicircum} conform, but no test of eco checks them.\\
\hline
\end{tabular}
\end{center}

% ===========================================================================
\section{Equality}\label{sec:appendix-implementation-equality}

The body requires that equality be structural, with IEEE equality at Float;
the identity of values plays no part, so a value containing NaN is not equal
to itself. When an equality test reaches two functions, its result is an
unspecified Bool: $\APPLY$ gives the set of both Bools
(Section~\ref{sec:dynamic-core-basic}, rule~\ruleref{r:d-app-basic-apply}
and Section~\ref{sec:appendix-dynamic-basis-apply}). % A50, A51

\subsection{Identity}
% A50
\begin{center}\small
\begin{tabular}{@{}>{\raggedright\arraybackslash}p{0.10\textwidth}>{\raggedright\arraybackslash}p{0.43\textwidth}>{\raggedright\arraybackslash}p{0.39\textwidth}@{}}
\hline
Impl. & Choice & Consequences for programs\\
\hline
Elm/JS & Structural, but \kw{True} at once when both operands are the same
object. At Float, \kw{x == 0} with \kw{0} written as an Int constant tests
whether \kw{x} is zero or NaN. &
Does not conform: \kw{p == p} with \kw{p = ( nan, 1 )} is \kw{True}, and
\kw{nan == 0} is \kw{True} (\kw{nan /= 0} is \kw{False}).\\
eco & Structural, but its code gives \kw{True} at once when both operands
are the same object, in the equality kernel and in the inline fast path;
whether two occurrences are the same object depends on optimisation. IEEE at
Float on every path. &
Does not conform for values containing NaN: by the code, \kw{p == p} as above
is \kw{True}, while a separately built \kw{( nan, 1 )} is unequal to \kw{p}.
Two of eco's tests of known Elm bugs expect \kw{False} for \kw{p == p} and for
other containers compared with themselves, as the definition gives (A50), but
eco's test harness does not run them.\\
% eco tests: elm-bugs CompilerIssue2362NanTupleEqualityTest:13-15 and
% CoreIssue988NanContainerSelfEqualityTest:17-24, not run by the E2E harness
% (survey, Equality and test suite status); code: Utils.cpp:503,
% EcoBackend.cpp:2216-2236. Mechanisation F7.
\hline
\end{tabular}
\end{center}

\subsection{Functions}
% A51
\begin{center}\small
\begin{tabular}{@{}>{\raggedright\arraybackslash}p{0.10\textwidth}>{\raggedright\arraybackslash}p{0.43\textwidth}>{\raggedright\arraybackslash}p{0.39\textwidth}@{}}
\hline
Impl. & Choice & Consequences for programs\\
\hline
Elm/JS & \kw{True} for the same object; otherwise the program stops with the
message ``Trying to use \kw{(==)} on functions''. &
Does not conform: \kw{(\textbackslash x -> x) == (\textbackslash x -> x + 0)}
and \kw{Just f == Just g} end the program.\\
eco & \kw{True} if the two are the same object, else \kw{False}, at any depth.
A lambda that captures no variable is one object however often it is
evaluated; a partial application is a new object each time. &
Conforms: the result is one member of the set. \kw{f == f} is \kw{True};
\kw{(\textbackslash x -> x) == (\textbackslash x -> x + 0)} is \kw{False}.\\
\hline
\end{tabular}
\end{center}

% ===========================================================================
\section{Bitwise Operations and String Conversion}%
\label{sec:appendix-implementation-bitwise}

The functions of the module \kw{Bitwise} are outside the initial basis
(Chapter~\ref{ch:appendix-static-basis}), and so are the conversions
\kw{String.toInt}, \kw{String.fromInt} and \kw{String.fromFloat}; the body
places no requirement on them. % A55, A64
They are recorded here because they reveal the Int range.

\subsection{Bitwise Operations}
% A55
\begin{center}\small
\begin{tabular}{@{}>{\raggedright\arraybackslash}p{0.10\textwidth}>{\raggedright\arraybackslash}p{0.43\textwidth}>{\raggedright\arraybackslash}p{0.39\textwidth}@{}}
\hline
Impl. & Choice & Consequences for programs\\
\hline
Elm/JS & 32 bits: operands converted to 32-bit integers, results signed
32-bit, shift counts taken modulo 32; \kw{shiftRightZfBy} gives an unsigned
32-bit result. &
\kw{shiftLeftBy 32 1} is \kw{1}; \kw{shiftRightZfBy 0 -1} is
\kw{4294967295}.\\
eco & 64 bits, signed; shift counts taken modulo 64; \kw{shiftRightZfBy} gives
its 64-bit result read as signed. &
\kw{shiftLeftBy 32 1} is \kw{4294967296}, \kw{shiftLeftBy 63 1} is the least
Int, and \kw{shiftRightZfBy 0 -1} is \kw{-1}.\\
\hline
\end{tabular}
\end{center}

\subsection{Conversion between Numbers and Strings}
% A64
\begin{center}\small
\begin{tabular}{@{}>{\raggedright\arraybackslash}p{0.10\textwidth}>{\raggedright\arraybackslash}p{0.43\textwidth}>{\raggedright\arraybackslash}p{0.39\textwidth}@{}}
\hline
Impl. & Choice & Consequences for programs\\
\hline
Elm/JS & \kw{String.toInt} reads an optional sign and decimal digits into a
double, with no overflow check and no length limit; \kw{"-0"} gives a negative
zero. \kw{String.fromInt} and \kw{String.fromFloat} use the host's printing of
numbers. &
Long numerals give rounded results rather than \kw{Nothing}.
\kw{String.fromInt} prints $10^{21}$ as \kw{"1e+21"}, and an Int holding
\kw{0.5} as \kw{"0.5"}.\\
eco & \kw{String.toInt} gives \kw{Nothing} when the numeral is outside
$\IntRange$ or longer than 64 code units, and the exact value otherwise.
\kw{String.fromInt} prints the exact decimal. \kw{String.fromFloat} prints the
shortest digits that read back to the same Float: fixed notation for
magnitudes from $10^{-6}$ up to $10^{21}$, exponent notation otherwise, and
\kw{"0"} for $-0$. &
\kw{String.toInt} and \kw{String.fromInt} are exact over $\IntRange$.\\
\hline
\end{tabular}
\end{center}

% ===========================================================================
\section{Evaluation Order and Debug}\label{sec:appendix-implementation-order}

The body requires the following (Section~\ref{sec:dynamic-core-rules},
Section~\ref{sec:dynamic-core-basic} and Chapter~\ref{ch:dynamic-modules}).
\begin{itemize}
\item The order in which the parts of a phrase are evaluated carries no
  meaning; the groups of a \kw{let} are evaluated in an order consistent with
  their dependencies, and the order among independent bindings is not
  specified. % A57, A58
\item Each top-level value is evaluated once, when its module is initialised,
  in dependency order; an implementation may evaluate it later or repeatedly
  provided the results agree. % A59
\item \kw{Debug.log l v} evaluates to $v$, and any output it makes is outside
  the definition; \kw{Debug.todo s} has no value; \kw{Debug.toString} gives an
  implementation-dependent String. % A62
\end{itemize}
Since evaluation has no other effect, the choices below are observable only
through the output of \kw{Debug.log}, through which of several subexpressions
without a value is reached first, and through resources.

\subsection{The Order of Evaluation within a Phrase}
% A57, A58
\begin{center}\small
\begin{tabular}{@{}>{\raggedright\arraybackslash}p{0.10\textwidth}>{\raggedright\arraybackslash}p{0.43\textwidth}>{\raggedright\arraybackslash}p{0.39\textwidth}@{}}
\hline
Impl. & Choice & Consequences for programs\\
\hline
Elm/JS & \kw{let} bindings by dependency groups, dependencies first. Record
fields by field name; tuple and list components, and the function and
arguments of an application, left to right. &
Conforms.\\
eco & \kw{let} bindings by dependency groups, dependencies first; named
bindings before bindings \kw{\_ = e}, and not in source order among
themselves; a run of bindings \kw{\_ = e} in source order. Record fields of
type Int, Float and Char first, then the rest, each by field name. Tuple and
list components and arguments in no fixed order. The callback of a basis
function such as \kw{List.map} may be applied in another order when it makes
no use of \kw{Debug}. &
Conforms. The output of several \kw{Debug.log} calls in one \kw{let} may
appear in a different order from Elm/JS.\\
\hline
\end{tabular}
\end{center}

\subsection{Top-level Values}
% A59, A25
\begin{center}\small
\begin{tabular}{@{}>{\raggedright\arraybackslash}p{0.10\textwidth}>{\raggedright\arraybackslash}p{0.43\textwidth}>{\raggedright\arraybackslash}p{0.39\textwidth}@{}}
\hline
Impl. & Choice & Consequences for programs\\
\hline
Elm/JS & Each once, at start-up, before \kw{main} runs, in dependency order. &
Conforms. A top-level value whose evaluation reaches \kw{Debug.todo} stops the
program at start-up, used or not.\\
eco & On first use. A value of type Int, Float or Char is computed again at
each use; a top-level value of function type evaluates its body again at each
call; any other value is computed once. &
Conforms. A \kw{Debug.log} in a top-level value may print once per use instead
of once.\\
\hline
\end{tabular}
\end{center}
In both, a \kw{main} in a module other than the main module is evaluated once,
and a top-level value needed while it is itself being evaluated, which has no
value, ends the program: by exhausting the stack, or, in Elm/JS compiled for
debugging, with a report of infinite recursion. % A25

\subsection{Debug}
% A62
\begin{center}\small
\begin{tabular}{@{}>{\raggedright\arraybackslash}p{0.10\textwidth}>{\raggedright\arraybackslash}p{0.43\textwidth}>{\raggedright\arraybackslash}p{0.39\textwidth}@{}}
\hline
Impl. & Choice & Consequences for programs\\
\hline
Elm/JS & \kw{Debug.log l v} writes \kw{l: } followed by \kw{Debug.toString v}
to the host's console. \kw{Debug.todo s} stops the program with ``TODO in
module \ldots{} on line \ldots'' and \kw{s}. A program using \kw{Debug} is
rejected when compiled with \kw{-{}-optimize}. &
Conforms. A program using \kw{Debug} cannot be compiled for production.\\
eco & \kw{Debug.log l v} writes \kw{l: } followed by \kw{Debug.toString v} to
the program's output when it is evaluated. \kw{Debug.todo s} writes
``Debug.todo: '' and \kw{s} to the error output and ends the process with
status 1. \kw{Debug} is accepted with \kw{-{}-optimize}, and its output
kept. &
Conforms.\\
\hline
\end{tabular}
\end{center}
\kw{Debug.toString} gives the same text in both for the forms tested, such as
\kw{\kwlb{} age = 3, name = "x" \kwrb}, \kw{Just -1} and \kw{0} for $-0$; for
Floats it follows \kw{String.fromFloat}
(Section~\ref{sec:appendix-implementation-bitwise}).

% ===========================================================================
\section{Resource Limits}\label{sec:appendix-implementation-limits}

The body places no bound on the size of a phrase
(Section~\ref{sec:syntax-core-restrictions}) or on the resources evaluation
uses; an implementation that stops for lack of resources is outside the
definition (Section~\ref{sec:dynamic-core-rules}). % A13, A34, A61

\subsection{The Size of Phrases}
% A13, A34, A39
\begin{center}\small
\begin{tabular}{@{}>{\raggedright\arraybackslash}p{0.10\textwidth}>{\raggedright\arraybackslash}p{0.43\textwidth}>{\raggedright\arraybackslash}p{0.39\textwidth}@{}}
\hline
Impl. & Choice & Consequences for programs\\
\hline
Elm/JS & No limit. & \\
eco & On every target, a program is rejected that has a function or lambda
with more than 2047 parameters, a lambda whose parameters and captured
variables together exceed 2047, a record type or record with more than 2047
fields, a constructor with more than 2040 fields, or a nullary constructor
that is not among the first 1024 constructors of its type. For native code
the closure limit is checked again after each polymorphic function is copied
for every type at which it is used. &
Programs beyond these limits are rejected. A closure within the limit in the
source can exceed it in native code: 700 local polymorphic functions each used
at three types inside one lambda make 2100 captured variables.\\
\hline
\end{tabular}
\end{center}

\subsection{Stack and Memory}
% A61
\begin{center}\small
\begin{tabular}{@{}>{\raggedright\arraybackslash}p{0.10\textwidth}>{\raggedright\arraybackslash}p{0.43\textwidth}>{\raggedright\arraybackslash}p{0.39\textwidth}@{}}
\hline
Impl. & Choice & Consequences for programs\\
\hline
Elm/JS & The host's stack and memory. Deep recursion that is not a self tail
call stops with a \kw{RangeError} the host can catch, at about ten thousand
frames. A \kw{Cmd.batch} nested 100000 deep exhausts the stack. &
Recursion over a long list must be a self tail call, which runs in constant
stack.\\
eco & The program runs on a stack of 64\,MiB. Exhausting it ends the process
with a segmentation fault; exhausting memory ends it with the message
``eco: out of memory'' and an abort. Deeply nested \kw{Cmd.batch} values are
handled without recursion. &
Far deeper recursion succeeds than in Elm/JS, but its failure cannot be
caught. Self tail calls run in constant stack.\\
\hline
\end{tabular}
\end{center}

% ===========================================================================
\section{Platform and Ports}\label{sec:appendix-implementation-platform}

The body defines tasks, commands, subscriptions and the encoding of values at
ports in Chapter~\ref{ch:effects}, and the execution of a program in
Section~\ref{sec:programs-execution}. It requires only that messages be
handled one at a time, each update seeing the model the previous one
produced, and the effects of an earlier update being dispatched before those
of a later one. % A86
It leaves the following unspecified or weakly specified.
\begin{itemize}
\item The order in which the commands of one batch are performed, and the
  relative order of effects handled by different managers. % A66, A67
\item An incoming port value is an event, handled after the step in progress
  completes, never during it. % A68
\item Outgoing port values are delivered after the step that produced them,
  in order per port; what the host does with them is outside the definition.
  % A69, A70
\item An incoming value that does not match the port's type is rejected at
  the boundary; the program is unaffected. % A71
\item Absent flags are the JSON value null. % A72
\item \kw{Bytes} is a payload type in every position, in ports and in flags,
  and its values cross the boundary as byte sequences, which are not JSON
  values; how a host represents a byte sequence is outside the definition
  (Section~\ref{sec:effects-ports}). % A92
\item JSON numbers are finite decimal rationals. A Float that is NaN or an
  infinity is encoded as \kw{null}, and $-0$ as the number $0$; an incoming
  number decodes at \kw{Int} exactly when it is an integer in $\IntRange$, an
  incoming zero being the one Int zero, and at \kw{Float} as the nearest float,
  $+0$ for a zero and an infinity for a number beyond the finite floats
  (Section~\ref{sec:effects-ports}). % A73, A74, A88
\item JSON objects are unordered, and error messages are unspecified. % A76
\item A spawned process that fails ends, and nothing else happens; killing a
  finished process has no effect; the task \kw{Process.sleep d} waits at least
  $\max(0,d)$ milliseconds, and for non-finite $d$ the delay is unspecified.
  % A79, A80, A81
\item The set of managers, beyond the task manager and ports, is a
  parameter of the platform. % A83
\end{itemize}

\subsection{The Order of Commands}
% A66, A67
\begin{center}\small
\begin{tabular}{@{}>{\raggedright\arraybackslash}p{0.10\textwidth}>{\raggedright\arraybackslash}p{0.43\textwidth}>{\raggedright\arraybackslash}p{0.39\textwidth}@{}}
\hline
Impl. & Choice & Consequences for programs\\
\hline
Elm/JS & A manager receives the commands of a batch in reverse order:
\kw{Cmd.batch [ perform a, perform b ]} runs \kw{b} first, and
\kw{Cmd.batch [ out 1, out 2 ]} delivers \kw{2} then \kw{1}. Several
subscriptions to one incoming port receive a value last first. Managers are
served in the order the program declares them and its ports. &
Conforms. Programs that depend on these orders, or on the order of port
declarations, behave differently in eco.\\
eco & The same reverse order within a batch and among subscriptions to one
port. Managers are served in no fixed order. &
Conforms.\\
\hline
\end{tabular}
\end{center}

\subsection{Incoming Port Values}
% A68, A71
\begin{center}\small
\begin{tabular}{@{}>{\raggedright\arraybackslash}p{0.10\textwidth}>{\raggedright\arraybackslash}p{0.43\textwidth}>{\raggedright\arraybackslash}p{0.39\textwidth}@{}}
\hline
Impl. & Choice & Consequences for programs\\
\hline
Elm/JS & The host's \kw{send} decodes the value and runs \kw{update} for each
current subscription before it returns. A value that does not decode makes
\kw{send} throw ``Trying to send an unexpected type of value through port''
to the host; the program is unaffected. &
Does not conform for delivery: a value sent from inside an outgoing
subscriber is handled during the step in progress, and a reply can be lost
because the subscription for it is not yet installed (elm/core issues 1110,
1141, 1154). Conforms for mismatches.\\
eco & The value is checked when the host sends it, then queued, and delivered
after the work in progress; a send from inside a subscriber is queued too. A
mismatch is reported to the host by an error code and the same message, and
the program is unaffected; but when the port's type cannot be checked at the
send, the value is checked on delivery, and a mismatch then ends the
process. &
Conforms for delivery. Does not conform for a mismatch found on delivery.\\
\hline
\end{tabular}
\end{center}

\subsection{Outgoing Port Values}
% A69, A70
\begin{center}\small
\begin{tabular}{@{}>{\raggedright\arraybackslash}p{0.10\textwidth}>{\raggedright\arraybackslash}p{0.43\textwidth}>{\raggedright\arraybackslash}p{0.39\textwidth}@{}}
\hline
Impl. & Choice & Consequences for programs\\
\hline
Elm/JS & Delivered later: a port waits one timer tick before its first
delivery, so that the host can subscribe after the program starts, and after
each batch. A value sent when the port has no subscriber is discarded. &
Conforms.\\
eco & Delivered at once, during the dispatch of the update that issued it; the
program waits while subscribers run. Values are kept until the port's first
subscriber and then handed to it in order; once a port has had a subscriber,
a value sent when it has none is discarded. &
Conforms.\\
\hline
\end{tabular}
\end{center}

\subsection{Values at the Boundary}
% A73, A74, A76, A88
\begin{center}\small
\begin{tabular}{@{}>{\raggedright\arraybackslash}p{0.10\textwidth}>{\raggedright\arraybackslash}p{0.43\textwidth}>{\raggedright\arraybackslash}p{0.39\textwidth}@{}}
\hline
Impl. & Choice & Consequences for programs\\
\hline
Elm/JS & Host values are JavaScript values. An incoming Int is any finite
number with no fractional part. A \kw{Json.Encode.Value} passes any host value
through, functions and \kw{undefined} included. Objects are written with
their keys in field-name order; of several missing fields, the last by name
is reported. &
Outgoing Ints above $2^{53}$ are inexact. NaN and the infinities pass as
Floats in both directions. Does not conform: an outgoing NaN or infinity is
delivered as a host number, where the definition delivers \kw{null}; an
outgoing $-0$ is delivered as $-0$; an incoming $-0$ at an \kw{Int} port
gives a negative zero, and at a \kw{Float} port $-0$, where the definition
gives $0$ and $+0$. An incoming NaN or infinity is a host value outside JSON,
and outside the definition.\\ % A88
eco & Host values are JSON text, byte sequences apart
(Section~\ref{sec:appendix-implementation-bytes}). An incoming Int is an integer, or a float
with an integral value, in $\IntRange$. An outgoing NaN or infinity is written
as \kw{null}. Keys in field-name order; missing fields are reported in the
order of the type's fields. &
\kw{9223372036854775808} and \kw{1e300} are rejected as Ints. A host value
that is not JSON cannot be sent. Conforms for NaN and the infinities.\\
\hline
\end{tabular}
\end{center}

\subsection{Bytes at the Boundary}\label{sec:appendix-implementation-bytes}
% A92
The body admits \kw{Bytes} wherever a payload type may appear, alone or
inside tuples, closed records, \kw{Maybe}, \kw{List} and \kw{Array}, in
the types of ports and in the flags of \kw{main}
(Section~\ref{sec:static-modules-ports}); its values cross as byte sequences
(Section~\ref{sec:effects-ports}). This rule follows the implementations that
provide binary data at the boundary. Each implementation chooses the host
representation of a byte sequence; in the table, \emph{eco/JS} is eco's
output for JavaScript.
\begin{center}\small
\begin{tabular}{@{}>{\raggedright\arraybackslash}p{0.10\textwidth}>{\raggedright\arraybackslash}p{0.43\textwidth}>{\raggedright\arraybackslash}p{0.39\textwidth}@{}}
\hline
Impl. & Choice & Consequences for programs\\
\hline
Elm/JS & No representation: a port or flags type in which \kw{Bytes}
occurs is rejected. &
Does not conform: every program with \kw{Bytes} in a payload or flags type
is rejected.\\
eco/JS & \kw{Bytes} in any payload position, in ports and in
flags; a byte sequence is a \kw{Uint8Array} at the host. &
Conforms.\\
eco & \kw{Bytes} only as the whole payload of a port, a byte sequence being
raw bytes at the host; \kw{Bytes} inside a payload type, and in flags, is
rejected. &
Does not conform: programs with \kw{Bytes} nested in a payload type, or in
flags, are rejected.\\
\hline
\end{tabular}
\end{center}

\subsection{Flags}
% A72
\begin{center}\small
\begin{tabular}{@{}>{\raggedright\arraybackslash}p{0.10\textwidth}>{\raggedright\arraybackslash}p{0.43\textwidth}>{\raggedright\arraybackslash}p{0.39\textwidth}@{}}
\hline
Impl. & Choice & Consequences for programs\\
\hline
Elm/JS & Absent flags are the JavaScript value \kw{undefined}, which no
decoder for \kw{null} or \kw{Maybe} accepts. Flags that do not decode make
\kw{init} throw ``Problem with the flags given to your Elm program on
initialization.'' &
Does not conform: a program with flags of type \kw{Maybe Int} started without
flags does not start.\\
eco & Absent flags are \kw{null}. Flags that do not decode stop an executable
with a message and an abort; an embedding host receives an error code and may
try again. &
Conforms: the program above starts with \kw{Nothing}.\\
\hline
\end{tabular}
\end{center}

\subsection{Processes}
% A79, A80, A81
\begin{center}\small
\begin{tabular}{@{}>{\raggedright\arraybackslash}p{0.10\textwidth}>{\raggedright\arraybackslash}p{0.43\textwidth}>{\raggedright\arraybackslash}p{0.39\textwidth}@{}}
\hline
Impl. & Choice & Consequences for programs\\
\hline
Elm/JS & A failing spawned process ends silently. Killing a finished or
already killed process throws a \kw{TypeError}. \kw{Process.sleep} uses the
host's timers: a delay of $0$, a negative delay, a non-finite delay or one
above $2^{31}-1$ ms is about 1 ms. &
Killing twice ends the program (elm/core issue 955): does not conform. A delay
above $2^{31}-1$ ms, about 24.8 days, is cut short: does not conform.\\
eco & A failing spawned process ends with a warning on the error output.
Killing a finished process has no effect. A delay of $0$ or less resumes at
once; a non-finite delay, or one above $2^{31}-1$ ms, is 1 ms. &
The warning is outside the definition: conforms. A delay above $2^{31}-1$ ms
is cut short: does not conform.\\
\hline
\end{tabular}
\end{center}
In both, a process runs until it waits for a timer, a message or another
primitive task, and is never suspended at an \kw{andThen} alone; this is one
of the behaviours the body allows. % A78

\subsection{Managers and Programs on a Platform}
% A83, A82
\begin{center}\small
\begin{tabular}{@{}>{\raggedright\arraybackslash}p{0.10\textwidth}>{\raggedright\arraybackslash}p{0.43\textwidth}>{\raggedright\arraybackslash}p{0.39\textwidth}@{}}
\hline
Impl. & Choice & Consequences for programs\\
\hline
Elm/JS & Every effect module the program uses is a manager: among them
\kw{Task}, \kw{Time}, \kw{Http}, \kw{Random}, \kw{Browser.Events} and
\kw{Browser.AnimationManager}. Browser programs run their \kw{init} and draw
their view; a \kw{main} of type \kw{Html msg} is drawn once. Several programs
may be started on one page. &
Every command and subscription of these managers is performed.\\
eco & In native code the managers are \kw{Task}, \kw{Time}, \kw{Http}, those
of eco's system package, and ports. A command or subscription for any other
manager, such as \kw{Random.generate} or \kw{Browser.Events}, is discarded.
\kw{Browser.element}, \kw{document} and \kw{application} return their record
without running \kw{init}; \kw{Browser.Navigation} ends the program with
status 1; \kw{Browser.Dom} tasks fail with \kw{NotFound}. A \kw{main} of type
\kw{Html msg} is evaluated and not drawn. One program per process. &
Only worker programs, \kw{Platform.worker}, run natively. Programs relying on
\kw{Random} or browser managers do nothing observable in those effects.\\
\hline
\end{tabular}
\end{center}

% ===========================================================================
\section{Extensions of eco}\label{sec:appendix-implementation-extensions}

The body admits no flags type beyond the port payload types
(Section~\ref{sec:static-modules-ports}); \kw{Bytes} is one of them, and is
not an extension (Section~\ref{sec:appendix-implementation-bytes}). % A33, A92
eco accepts the following programs, which are not Standard Elm.

\begin{center}\small
\begin{tabular}{@{}>{\raggedright\arraybackslash}p{0.10\textwidth}>{\raggedright\arraybackslash}p{0.43\textwidth}>{\raggedright\arraybackslash}p{0.39\textwidth}@{}}
\hline
Impl. & Choice & Consequences for programs\\
\hline
Elm/JS & None. & \\
eco & A program of type
\kw{Program System.Environment model msg}, with \kw{System.Environment} from
eco's system package, has its flags decoded by a built-in decoder, with no
check of the payload type. &
Programs using this extension are eco programs, not Standard Elm, and are
rejected by Elm 0.19.1.\\
\hline
\end{tabular}
\end{center}

% ===========================================================================
\section{Known Divergences}\label{sec:appendix-implementation-divergences}

The tables below list every known point on which an implementation does not
conform: the implementation (Elm/JS, eco, or both), the construct, what the
implementation does, what the definition says, and the ruling that settles
the point. Where a ruling chose a meaning neither implementation has, both
are listed.

% The divergence tables are split by area so that each fits a page.

\subsection{Int Arithmetic}
\begin{center}\footnotesize\setlength{\tabcolsep}{3pt}
\begin{tabular}{@{}>{\raggedright\arraybackslash}p{0.09\textwidth}>{\raggedright\arraybackslash}p{0.19\textwidth}>{\raggedright\arraybackslash}p{0.27\textwidth}>{\raggedright\arraybackslash}p{0.26\textwidth}>{\raggedright\arraybackslash}p{0.09\textwidth}@{}}
\hline
Impl. & Construct & Implementation & Definition & Ruling\\
\hline
Elm/JS & Int \kw{+ - *} beyond $\pm 2^{53}$ & rounded to a double; an
infinity beyond the largest double & $\norm$ of the exact result & A40\\
Elm/JS & integer constant at or above $2^{63}$ & reduced modulo $2^{64}$, then
held as a double & $\norm(n)$ & A9\\
Elm/JS & \kw{a // b}, quotient outside $\pm 2^{31}$ & reduced to 32 bits &
$\norm$ of the truncated quotient & A41\\
Elm/JS & \kw{modBy 0 x} & program stops & $0$ & A42\\
Elm/JS & \kw{remainderBy 0 x} & NaN held as an Int & $0$ & A42\\
Elm/JS & Int \kw{b \textasciicircum{} e}, $e<0$ & a fraction or an infinity &
$1$, $\pm 1$ or $0$ & A43\\
eco & Int \kw{b \textasciicircum{} e}, $e<0$, $b=\pm 1$ & $0$ & $1$, or $\pm 1$
by parity & A43\\
Elm/JS & \kw{round}, \kw{floor}, \kw{ceiling} of NaN or an infinity & NaN or
the infinity held as an Int & $0$, or a bound of $\IntRange$ & A44\\
Elm/JS & \kw{truncate} outside $\pm 2^{31}$ & reduced to 32 bits & saturated
to $\IntRange$ & A44\\
Elm/JS & \kw{negate 0} and other Int zeros & a negative zero & one zero &
A45\\
\hline
\end{tabular}
\end{center}

\subsection{Float, Order and Equality}
\begin{center}\footnotesize\setlength{\tabcolsep}{3pt}
\begin{tabular}{@{}>{\raggedright\arraybackslash}p{0.09\textwidth}>{\raggedright\arraybackslash}p{0.19\textwidth}>{\raggedright\arraybackslash}p{0.27\textwidth}>{\raggedright\arraybackslash}p{0.26\textwidth}>{\raggedright\arraybackslash}p{0.09\textwidth}@{}}
\hline
Impl. & Construct & Implementation & Definition & Ruling\\
\hline
Elm/JS & \kw{compare} with NaN & always \kw{GT} & total order, NaN greatest &
A46\\
Elm/JS & \kw{<}, \kw{<=}, \kw{>}, \kw{>=} with NaN & depends on whether an
operand is a constant & IEEE predicates & A47\\
Elm/JS & Float \kw{min}, \kw{max} with NaN or two zeros & depends on the
order of the arguments: \kw{min nan 1} and \kw{max 1 nan} are \kw{1} &
IEEE 754-2019 minimum and maximum: NaN for both & A48\\
Elm/JS & Float \kw{clamp} with $lo>hi$, a NaN argument, or zeros of
different signs & elm/core definition: \kw{clamp 10 0 5} is \kw{10},
\kw{clamp 0 1 nan} is \kw{1}, \kw{clamp nan 1 0.5} is \kw{0.5},
\kw{clamp 0 1 -0} is $-0$ & \kw{min (max x lo) hi} with IEEE 754-2019
minimum and maximum: $+0$, NaN, NaN, $+0$ & A47, A48, A49\\
Elm/JS & Float \kw{1 \textasciicircum{} (1/0)} & NaN & $1$ & A65\\
both & \kw{v == v} with $v$ containing NaN & \kw{True} when both operands are
one object & \kw{False} & A50\\
Elm/JS & Float \kw{x == 0}, \kw{0} an Int constant & \kw{True} for NaN &
\kw{False} for NaN & A50\\
Elm/JS & \kw{==} reaching two functions & program stops & an unspecified
Bool & A51\\
\hline
\end{tabular}
\end{center}

\subsection{Characters and Strings}
\begin{center}\footnotesize\setlength{\tabcolsep}{3pt}
\begin{tabular}{@{}>{\raggedright\arraybackslash}p{0.09\textwidth}>{\raggedright\arraybackslash}p{0.19\textwidth}>{\raggedright\arraybackslash}p{0.27\textwidth}>{\raggedright\arraybackslash}p{0.26\textwidth}>{\raggedright\arraybackslash}p{0.09\textwidth}@{}}
\hline
Impl. & Construct & Implementation & Definition & Ruling\\
\hline
both & \kw{\textbackslash u\kwlb D800\kwrb} and other surrogate escapes &
accepted & not an escape sequence & A8\\
Elm/JS & \kw{toCode} of a lone surrogate; Chars from case mapping & NaN; a
Char of two characters & \kw{toCode} injective into the integers & A52\\
Elm/JS & Char order & by UTF-16 code units & by \kw{toCode} & A53\\
both & String order & by UTF-16 code units & lexicographic by characters
ordered by \kw{toCode} & A53\\
\hline
\end{tabular}
\end{center}

\subsection{Evaluation}
\begin{center}\footnotesize\setlength{\tabcolsep}{3pt}
\begin{tabular}{@{}>{\raggedright\arraybackslash}p{0.09\textwidth}>{\raggedright\arraybackslash}p{0.19\textwidth}>{\raggedright\arraybackslash}p{0.27\textwidth}>{\raggedright\arraybackslash}p{0.26\textwidth}>{\raggedright\arraybackslash}p{0.09\textwidth}@{}}
\hline
Impl. & Construct & Implementation & Definition & Ruling\\
\hline
Elm/JS & \kw{(\&\&) a b}, \kw{(||) a b} applied in prefix form & second
argument evaluated only when needed: \kw{(\&\&) False (Debug.todo "x")} is
\kw{False} & ordinary strict functions: no value & A56\\
Elm/JS & lambdas made in a self tail-recursive function & share the
function's changing parameters: three closures give \kw{[3,3,3]} & each holds
the values of its own call: \kw{[3,2,1]} & A63\\
both & a value whose scheme quantifies a \kw{number} variable, used only at
Float, such as \kw{x = if 2 \textasciicircum{} -1 == 0 then Debug.todo "x"
else 1} & Elm/JS evaluates it once, with its one kind of number; eco only at
the instances used. Either way \kw{x} is $1$ and the program runs & the default
copy, at Int, is also evaluated; there \kw{2 \textasciicircum{} -1} is $0$,
so it has no value and the program has no meaning & A87\\
\hline
\end{tabular}
\end{center}

\subsection{Programs and Ports}
\begin{center}\footnotesize\setlength{\tabcolsep}{3pt}
\begin{tabular}{@{}>{\raggedright\arraybackslash}p{0.09\textwidth}>{\raggedright\arraybackslash}p{0.19\textwidth}>{\raggedright\arraybackslash}p{0.27\textwidth}>{\raggedright\arraybackslash}p{0.26\textwidth}>{\raggedright\arraybackslash}p{0.09\textwidth}@{}}
\hline
Impl. & Construct & Implementation & Definition & Ruling\\
\hline
Elm/JS & incoming port value & \kw{update} runs inside the host's \kw{send} &
an event handled after the step in progress & A68\\
eco & incoming value not matching the port, found on delivery & the process
ends & rejected at the boundary; program unaffected & A71\\
Elm/JS & absent flags & \kw{undefined}: \kw{Maybe} flags rejected & \kw{null}
& A72\\
Elm/JS & outgoing Float NaN or infinity & passed as a host number &
\kw{null} & A88\\
Elm/JS & $-0$ across a port & passed as $-0$; an incoming $-0$ at \kw{Int}
is a negative zero & the number $0$; $0$ at \kw{Int}, $+0$ at \kw{Float} &
A88\\
Elm/JS & \kw{Process.kill} of a finished process & program stops & no effect
& A80\\
both & \kw{Process.sleep d}, $d>2^{31}-1$ & about 1 ms & at least $d$ ms &
A81\\
both & a \kw{main} of unsupported type outside the main module & rejected &
an ordinary value & A32, A82\\
both & two ports of one name, or a port named after a manager & accepted;
the program stops at start-up & rejected by elaboration & A84\\
Elm/JS & \kw{Bytes} in a port or flags type & rejected & a payload type,
crossing as a byte sequence & A92\\
eco & \kw{Bytes} inside a payload type, or in flags, native code & rejected
& a payload type, crossing as a byte sequence & A92\\
\hline
\end{tabular}
\end{center}

\subsection{Elaboration}
\begin{center}\footnotesize\setlength{\tabcolsep}{3pt}
\begin{tabular}{@{}>{\raggedright\arraybackslash}p{0.09\textwidth}>{\raggedright\arraybackslash}p{0.19\textwidth}>{\raggedright\arraybackslash}p{0.27\textwidth}>{\raggedright\arraybackslash}p{0.26\textwidth}>{\raggedright\arraybackslash}p{0.09\textwidth}@{}}
\hline
Impl. & Construct & Implementation & Definition & Ruling\\
\hline
both & a pattern \kw{( )} with a gap & accepted as unit & not a phrase & A1\\
eco & a rigid type variable escaping its annotation & the compiler fails with
an internal error & does not elaborate & A20\\
eco & a record extended through a row variable by a field it already has &
accepted, the outer field winning & does not elaborate & A23\\
Elm/JS & unannotated mutually recursive values used at several types &
accepted; some such programs fail at run time & each member monomorphic
within its group & A24\\
eco & two literal patterns that denote one value, spelt differently & second
not reported redundant & rejected as redundant & A31\\
eco & a constrained type variable as a row variable, \kw{\kwlb{} number | x
: Int \kwrb} & annotation accepted; every use fails & ill formed & A36\\
eco & polymorphic recursion through annotated functions, native code &
rejected & elaborates & A39\\
eco & a closure over 2047 variables after copying polymorphic functions,
native code & rejected & elaborates & A39\\
eco & phrases beyond the size limits of
Section~\ref{sec:appendix-implementation-limits} & rejected & elaborate & A13,
A34\\
both & \kw{import} \mv{modid} \kw{as Basics} for a module other than Basics
& accepted (not tested) & does not elaborate & A91\\
\hline
\end{tabular}
\end{center}
